\chapter{Marco teórico o Estado del arte}
\label{cap:estado}

Introducción al marco teórico, justificación de la elección de los temas tratados y su relación con el proyecto. Se trata de un capítulo de investigación documental que sirve como base para la toma de decisiones de diseño y desarrollo del prototipo. Consta de cuatro secciones principales: fundamentos de la narrativa digital y videojuegos, salud mental y perspectiva de género en los videojuegos, sistemas de juego dinámicos y mecánicas adaptativas, y análisis de obras de referencia. Todas ellas se han seleccionado por su relevancia en relación con la problemática planteada y la solución propuesta, así como por su capacidad para proporcionar un marco conceptual sólido que guíe el diseño y desarrollo del prototipo.

\section{Fundamentos de la Narrativa Digital y Videojuegos}
Con la llegada de los medios interactivos, la estructura narrativa ha evolucionado hacia el concepto de la Narrativa Digital (\textit{Digital Storytelling}), definida como la práctica de combinar la narración con contenidos digitales (imágenes, sonido, vídeo e interactividad) para construir una obra con un fuerte componente emocional y una finalidad que a menudo va más allá del simple entretenimiento: concienciar, educar o generar empatía \parencite{navarro2023}. A diferencia de los medios tradicionales cuya lógica se fundamenta en la representación pasiva, los videojuegos permiten una interacción activa del jugador, integrando al usuario como un agente causal que ejecuta y habita el relato en tiempo real \parencite{gomez2026, hernandez2024}.

A la capacidad de los videojuegos para transmitir argumentos y persuadir mediante sus reglas y sistemas se la conoce como retórica procedimental, concepto que ha sido analizado por \textcite{bogost2011}.

Existen tres pilares esenciales en toda estructura narrativa digital \parencite{navarro2023}:
\begin{itemize}
  \item \textbf{Contenido narrativo y modelo de mundo:} diferencia la cronología de vivencias del protagonista de la trama mostrada al jugador. El mundo del juego representa el conflicto íntimo del personaje, empleando un entorno apocalíptico como alegoría de sus dilemas.
  \item \textbf{Tiempo y ritmo (\textit{pacing}):} relación entre el tiempo de la historia y el tiempo del relato, en este caso del videojuego.
  \item \textbf{Personajes:} el jugador debe poder identificarse con los problemas, dudas y el agotamiento vital del protagonista.
\end{itemize}

\subsection{Ludología, Narratología y Armonía Ludonarrativa}
La ludología y la narratología son dos enfoques complementarios en el estudio de los videojuegos \parencite{finestrat2018, hernandez2024}. Tradicionalmente, la ludología situó el centro de la experiencia en el sistema de reglas, la simulación de comportamientos y la capacidad de ejercer el control, mientras que la narratología se enfocó en la estructura dramática, los personajes y la estética de la historia \parencite{cuadrado2020, gomez2026}.

De la fricción o el equilibrio entre ambas dimensiones surgen dos fenómenos clave estudiados por \textcite{hernandez2024}: la \textit{disonancia ludonarrativa}, que se produce cuando la estructura narrativa y las mecánicas jugables entran en contradicción rompiendo el significado de ello \parencite{gomez2026, hernandez2024}, y la \textit{consonancia ludonarrativa} (o armonía ludonarrativa), que surge cuando existe una alta cohesión entre la historia que se desea contar y las mecánicas e interfaz (incluyendo elementos diegéticos), dando a la obra más profundidad y significado artístico \parencite{hernandez2024}. Asimismo, tal y como demuestran \textcite{hernandez2024}, incluso estas fricciones ludonarrativas pueden emplearse deliberadamente como un recurso poético y expresivo para despertar emociones concretas en el jugador.

En este proyecto, se busca alcanzar una consonancia ludonarrativa profunda mediante la implementación de mecánicas adaptativas e interfaces que reflejen el estado psicológico del protagonista y las decisiones del jugador, permitiendo que la narrativa y la jugabilidad se complementen y potencien mutuamente \parencite{hernandez2024}.

\subsection{Narrativa Ambiental y Metáfora Visual}
La narrativa ambiental y la metáfora visual son técnicas utilizadas en los videojuegos para contar historias y transmitir emociones a través del diseño del entorno y los elementos visuales \parencite{finestrat2018, boyero2025}. Estas técnicas permiten a los jugadores experimentar la trama de manera más inmersiva, apoyándose en la abstracción gráfica, el cierre visual y la ambigüedad perceptiva para que el usuario descubra e interprete progresivamente el mundo del juego como símbolo del conflicto mental del protagonista \parencite{boyero2025}.

Obras como \textit{Inside} \parencite{playdead2016} y \textit{Little Nightmares} \parencite{tarsier2017} emplean la narrativa ambiental y la metáfora visual para explorar temas de opresión, control y miedo, utilizando el diseño de niveles, la iluminación y los elementos visuales para crear una atmósfera que refleje los dilemas del protagonista y genere empatía en el jugador \parencite{boyero2025}.

Para lograr una narrativa ambiental efectiva, es fundamental considerar aspectos como la coherencia del mundo del juego, la simbología de los elementos visuales y la interacción del jugador con el entorno \parencite{finestrat2018, boyero2025}. En este caso, el mundo del juego se presenta como un entorno apocalíptico de zombies donde tras recibir un mordisco el protagonista empieza a tener alucinaciones y da lugar a entornos abstractos mutantes, que refleja los conflictos internos del protagonista relacionados con su identidad de género y salud mental, convirtiendo el espacio transitable en una representación simbólica de sus luchas internas.

\subsection{Empatía, Agencia e Identificación con el Personaje}
El presente proyecto busca fomentar la empatía, la agencia y la identificación del jugador con el protagonista del juego. Como señala \textcite{cuadrado2020}, resulta necesario reformular el concepto clásico de actuación o agencia (\textit{agency}): más allá de entenderla únicamente como libertad absoluta de control mecánico (agencia funcional), la agencia en los videojuegos narrativos se abre hacia una agencia afectiva o dramática (\textit{dramatic agency}) vinculada a la construcción de sentido y al compromiso emocional con el personaje y su historia \parencite{cuadrado2020, hernandez2024}.

En esta línea, \textcite{gomez2026} propone que la construcción de la identidad en el videojuego se forma a través de una «tríada ontológica» compuesta por el \textit{Jugador} (entidad real), el \textit{Personaje} (entidad diegética con su trasfondo) y el \textit{Avatar} (cuerpo digital mediador). A través de esta relación triádica surge una «identidad proyectiva» en la que el usuario vierte sus valores y gestiona su propio «yo» mediante la toma de decisiones, llegando a generar un «residuo identitario», que consiste en un sedimento cognitivo, empático y moral que perdura en la conciencia del jugador tras finalizar la sesión de juego \parencite{gomez2026}.

Asimismo, como advierte \textcite{gomez2026} al analizar la «paradoja de la desposesión», cuando un sistema toma completo control del avatar al usuario, se produce una crisis en la agencia y una sensación de alienación. Es precisamente para evitar esta ruptura no deseada por lo que el presente \ac{TFG} introduce el mecanismo de \textit{Consentimiento Explícito}, asegurando que la transformación del avatar y sus mecánicas sea consensuada con el jugador.

\section{Salud Mental y Perspectiva de Género en los Videojuegos}
El análisis de los videojuegos desde la perspectiva de género ha cobrado una notable fuerza académica y creativa en la última década. Más allá de analizar la representación cuantitativa de personajes en la industria, aplicar la perspectiva de género al diseño de un videojuego implica cuestionar cómo las estructuras sociales, los roles de género normativos y la cisnormatividad impactan en la subjetividad y en la salud mental de las personas. Al operar como un «laboratorio ontológico» protegido dentro del círculo mágico lúdico, el videojuego permite al jugador habitar y experimentar de forma encarnada la fragmentación psicológica, la disforia y la presión normativa que atraviesa el protagonista \parencite{gomez2026}.

\section{Sistemas de Juego Dinámicos y Mecánicas Adaptativas}
En la historia de los videojuegos, se han desarrollado diversos sistemas de juego dinámicos y mecánicas adaptativas que permiten responder a las acciones del jugador de manera personalizada. Estos sistemas pueden ajustar la dificultad, modificar la narrativa o alterar la experiencia de juego en función del comportamiento del usuario, validando su identidad operativa a través de la retroalimentación sistémica (\textit{feedback}) \parencite{gomez2026, hernandez2024}.

En \textit{Disco Elysium} \parencite{zaum2020}, por ejemplo, las decisiones del jugador afectan directamente a la narrativa y fragmentan el diálogo psicológico interno del personaje \parencite{gomez2026}, mientras que en \textit{Silent Hill: Shattered Memories} \parencite{climax2009}, el juego adapta su atmósfera y los encuentros según las reacciones del jugador. Por otro lado, \textit{Undertale} \parencite{fox2015} ofrece un sistema que registra las decisiones mecánicas del jugador para interpelar su responsabilidad ética y alterar la historia, los diálogos y la interfaz en consecuencia \parencite{gomez2026}, mientras que \textit{Tell Me Why} \parencite{dontnod2020} utiliza la narrativa y el diseño de personajes para explorar temas de identidad de género y salud mental. Asimismo, títulos como \textit{Inside} \parencite{playdead2016} y \textit{Little Nightmares} \parencite{tarsier2017} emplean la narrativa ambiental y la metáfora visual para transmitir emociones y contar historias profundas sin necesidad de diálogos \parencite{boyero2025}.

Por lo tanto, la implementación de sistemas de juego dinámicos y mecánicas adaptativas no solo mejora la experiencia del jugador, sino que también permite abordar temas complejos y sensibles de manera más efectiva, fomentando la empatía y la comprensión a través de la interacción lúdica \parencite{bogost2011, cuadrado2020, gomez2026}. En este contexto, se propone un sistema de juego que observe las acciones del jugador y modifique dinámicamente la experiencia —incluyendo la interfaz de usuario (\ac{UI}) como extensión del flujo de conciencia del avatar \parencite{gomez2026, hernandez2024} y las mecánicas— para reflejar el estado psicológico del protagonista.

\section{Análisis de Obras de Referencia}
Para fundamentar tanto el diseño artístico como el diseño técnico del proyecto, se han estudiado diversas obras de referencia que comparten puntos de contacto con este \ac{TFG}:
\begin{itemize}
  \item \textit{Undertale} \parencite{fox2015}: un ejemplo de cómo el sistema emplea la retórica procedimental para registrar las decisiones mecánicas del jugador (atacar, perdonar, huir), deconstruyendo las convenciones del género para alterar profundamente la narrativa, los diálogos e incluso la propia interfaz de usuario y las reglas del combate \parencite{gomez2026}.
  \item \textit{Silent Hill: Shattered Memories} \parencite{climax2009}: referente técnico fundamental en el uso de un sistema de perfilado psicológico en segundo plano. El juego observa silenciosamente qué objetos mira el jugador, cómo reacciona ante los peligros y cómo explora el escenario, modificando dinámicamente la estética del mundo, el aspecto de las criaturas y el desenlace de la historia.
  \item \textit{Tell Me Why} \parencite{dontnod2020}: un juego de aventuras narrativas que aborda temas de identidad de género y salud mental, utilizando la narrativa y el diseño de personajes para explorar la experiencia de una persona transgénero.
  \item \textit{Inside} \parencite{playdead2016}: un juego de plataformas y aventuras que utiliza la narrativa ambiental y la metáfora visual para contar una historia profunda y emocional sin necesidad de diálogos, explorando temas de opresión y control \parencite{boyero2025}.
  \item \textit{Little Nightmares} \parencite{tarsier2017}: un juego de aventuras y plataformas que utiliza la narrativa ambiental y la metáfora visual para crear una atmósfera inquietante y transmitir emociones a través del diseño de niveles y la interacción con el entorno \parencite{boyero2025}.
  \item \textit{Disco Elysium} \parencite{zaum2020}: un juego de rol que combina narrativa profunda con mecánicas adaptativas que representan la fragmentación psicológica del protagonista, permitiendo al jugador influir en la historia y en la configuración interna del personaje a través de sus decisiones \parencite{gomez2026}.
\end{itemize}